\documentclass[11pt]{amsart}
\usepackage[T1]{fontenc}
\usepackage{lmodern,microtype,amsmath,amssymb,amsthm,mathtools,enumitem,booktabs}
\usepackage[margin=1.1in]{geometry}
\usepackage[hidelinks]{hyperref}
\usepackage{xurl}

\newtheorem{theorem}{Theorem}
\numberwithin{equation}{section}

\newcommand{\Pp}{\mathbb P}
\newcommand{\E}{\mathbb E}
\newcommand{\dbar}{\bar d}

\hypersetup{pdftitle={The variation threshold, Kemeny's constant and stationary times of reversible Markov chains},pdfauthor={Yitzchak Shmalo}}

\title[The variation threshold and stationary times]{The variation threshold, Kemeny's constant and stationary times of reversible Markov chains}
\author{Yitzchak Shmalo}
\address{Einstein Institute of Mathematics, The Hebrew University of Jerusalem, Jerusalem, Israel}
\email{yitzchak.shmalo@gmail.com}
\date{13 October 2026}
\subjclass[2020]{60J27, 60J10}
\keywords{Reversible Markov chain, mixing time, Kemeny's constant, relaxation time, stationary time}

\begin{document}
\begin{abstract}
For finite reversible Markov chains, Aldous and Fill asked whether the variation threshold $\tau_1$ is at most the average hitting time $\tau_0$ (Kemeny's constant), and whether the relaxation time $\tau_2$ is at most $\tau_1^{(2)}$, the largest over initial states of the least mean of a stopping time at which the chain has its stationary distribution. The answer to the first question is no: a birth and death chain on three states with a strong drift has $\tau_1>\tau_0$, and on the analogous chains with $k+1$ states the ratio $\tau_1/\tau_0$ tends to $t_k/k$ as the drift grows, where $t_k$ is the point at which the survival function of the Gamma$(k,1)$ distribution equals $e^{-1}$; for $k=4$ this limit is $1.0879\ldots$. The answer to the second question is yes, in continuous and in discrete time.
\end{abstract}
\maketitle

\section{Introduction}

Let $X$ be an irreducible Markov chain in continuous time on a finite set $I$ of $n\ge2$ states, with generator $Q=(q_{ij})$ and stationary distribution $\pi$, and assume that $X$ is reversible: $\pi_iq_{ij}=\pi_jq_{ji}$ for all $i,j$. Write $\Pp_i$ and $\E_i$ for the chain started at $i$, $T_j=\inf\{t\ge0:X_t=j\}$ for the hitting time of $j$, $\|\mu-\nu\|=\max_A|\mu(A)-\nu(A)|$ for the variation distance between probability measures on $I$, $\mu(h)=\sum_k\mu_kh(k)$, and $0=\lambda_1<\lambda_2\le\dots\le\lambda_n$ for the eigenvalues of $-Q$. In the notation of Aldous and Fill \cite[Chapter~4]{AF}, the average hitting time, the variation threshold and the relaxation time are
\[
\tau_0=\sum_{i,j}\pi_i\pi_j\E_iT_j,\qquad \tau_1=\min\{t\ge0:\dbar(t)\le e^{-1}\},\qquad \tau_2=\frac1{\lambda_2},
\]
where $\dbar(t)=\max_{i,j}\|\Pp_i(X_t\in\cdot)-\Pp_j(X_t\in\cdot)\|$. By the random target lemma \cite[Corollary~2.14]{AF}, $\sum_j\pi_j\E_iT_j$ does not depend on $i$, so $\tau_0$ is this common value, often called Kemeny's constant. We also need
\[
\tau_1^{(2)}=\max_i\,\inf_U\E_iU,
\]
where the infimum is over randomized stopping times $U$ with $\Pp_i(X_U\in\cdot)=\pi$. A randomized stopping time is a stopping time for a filtration $(\mathcal F_t)$ with respect to which $X$ is a Markov process, for instance the filtration generated by $X$ and a random variable independent of $X$.

Among the inequalities between these parameters proved in \cite[Chapter~4]{AF} are
\[
\tau_2\le\tau_1\le66\,\tau_0,\qquad\tau_2\le\tau_0,\qquad\tau_1^{(2)}\le\tau_0
\]
(Lemmas~4.23, 4.15 and~4.24 and Section~4.3.2 there), and by \cite[Theorem~4.6]{AF} the parameter $\tau_1^{(2)}$ belongs to a family of parameters equivalent to $\tau_1$, that is, bounded above and below by absolute constant multiples of it. Aldous and Fill chose $\tau_1$ as the representative of this family partly because it gives the constant $1$ in $\tau_2\le\tau_1$, and they ask whether the constants in other basic inequalities relating the family to other parameters can be made $1$ as well. Their Open Problem~4.48 asks:
\begin{enumerate}[label=(\alph*)]
\item Is $\tau_1\le\tau_0$?
\item Is $\tau_2\le\tau_1^{(2)}$?
\end{enumerate}
We show that the answer to (a) is no and the answer to (b) is yes.

For $k\ge1$ and $\varepsilon>0$ let $X^{(k,\varepsilon)}$ be the chain on $\{0,1,\dots,k\}$ that jumps from $i$ to $i-1$ at rate $1$ for $1\le i\le k$ and from $i$ to $i+1$ at rate $\varepsilon$ for $0\le i\le k-1$. It is reversible, with $\pi_i$ proportional to $\varepsilon^i$. Let $\Gamma_k$ be a random variable with the Gamma$(k,1)$ distribution, the law of a sum of $k$ independent standard exponential variables, and let $t_k$ be the solution of
\begin{equation}\label{eq:tk}
\Pp(\Gamma_k>t_k)=e^{-t_k}\sum_{j=0}^{k-1}\frac{t_k^j}{j!}=e^{-1}.
\end{equation}

\begin{theorem}\label{thm:a}
\begin{enumerate}[label=\textup{(\roman*)}]
\item If $0<\varepsilon\le1/100$, then $X^{(2,\varepsilon)}$ has
\[
\tau_0=\frac{2(1+\varepsilon)}{1+\varepsilon+\varepsilon^2}<2<\tau_1.
\]
\item For every $k\ge1$, the chain $X^{(k,\varepsilon)}$ has $\tau_0\to k$ and $\tau_1\to t_k$ as $\varepsilon\to0$.
\end{enumerate}
\end{theorem}

As $\varepsilon\to0$ the chain $X^{(k,\varepsilon)}$ becomes the pure death chain absorbed at $0$. For that chain $\dbar(t)=\Pp(\Gamma_k>t)$, the probability that the chain started at $k$ has not reached $0$ by time $t$, while $\tau_0$ tends to $k$, the mean of $\Gamma_k$; so $\tau_1>\tau_0$ for small $\varepsilon$ whenever $\Pp(\Gamma_k>k)>e^{-1}$, which holds for $k=2$ because $\Pp(\Gamma_2>2)=3e^{-2}=0.406\ldots$. For $\varepsilon=1/100$ the chain $X^{(2,\varepsilon)}$ has $\tau_1=2.1330\ldots$ and $\tau_1/\tau_0=1.0666\ldots$ (Section~\ref{sec:checks}). The limits $t_k/k$ in part~(ii) are $1$, $1.0730\ldots$, $1.0860\ldots$, $1.0879\ldots$ and $1.0868\ldots$ for $k=1,\dots,5$; they tend to $1$ as $k\to\infty$ by the law of large numbers, and among $k\le60$ the largest is $t_4/4=1.08799\ldots$. So the best constant $C$ such that $\tau_1\le C\tau_0$ for every finite reversible chain satisfies $t_4/4\le C\le66$. We do not know its value.

\begin{theorem}\label{thm:b}
Let $u$ be a real function on $I$ with $Qu=-\lambda_2u$ and $\max_k|u(k)|=1$. For every state $i$ and every randomized stopping time $U$ with $\Pp_i(X_U\in\cdot)=\pi$,
\[
\E_iU\ge\tau_2\,|u(i)|.
\]
In particular $\tau_2\le\tau_1^{(2)}$.
\end{theorem}

The constant $1$ cannot be lowered: two-state chains have $\tau_2=\tau_1^{(2)}$ (Remark~1 in Section~\ref{sec:b}). Since $\tau_1^{(2)}\le\tau_0$, Theorem~\ref{thm:b} refines the inequality $\tau_2\le\tau_0$.

Both theorems hold for chains in discrete time with the conventions of \cite[(3.39) and Section~4.3.1]{AF}. For an irreducible reversible transition matrix $P$ with second largest eigenvalue $\lambda_2$, the relaxation time is $\tau_2=1/(1-\lambda_2)$; $\tau_0$ and $\tau_1^{(2)}$ are defined as above, with hitting times and stopping times taking values in $\{0,1,2,\dots\}$; and $\tau_1$ is the variation threshold of the continuized chain, the chain in continuous time with generator $P-I$. Section~\ref{sec:a-disc} shows that the chains $I+cQ$, with $Q$ the generator of $X^{(k,\varepsilon)}$ and $0<c\le1/(1+\varepsilon)$, have the same ratio $\tau_1/\tau_0$ as $X^{(k,\varepsilon)}$, and that for $k=2$, $0<\varepsilon\le1/100$ and small $c$ the variation threshold of the discrete chain itself also exceeds $\tau_0$; Section~\ref{sec:b} proves Theorem~\ref{thm:b} in discrete time, with $Pu=\lambda_2u$ in place of $Qu=-\lambda_2u$.

\section{Proof of Theorem~\ref{thm:a}}\label{sec:a}

Write $Q_\varepsilon$ for the generator of $X^{(k,\varepsilon)}$ and $\dbar_\varepsilon$ for its function $\dbar$. Detailed balance, $\pi_i\varepsilon=\pi_{i+1}$, shows that $X^{(k,\varepsilon)}$ is reversible with $\pi_i\propto\varepsilon^i$. By the random target lemma and the eigentime identity \cite[Proposition~3.13]{AF},
\begin{equation}\label{eq:eigentime}
\tau_0=\sum_{m\ge2}\frac1{\lambda_m}.
\end{equation}

\subsection{The chain on three states}
Let $k=2$. The matrix $-Q_\varepsilon$ has trace $2+2\varepsilon$, principal $2\times2$ minors $\varepsilon^2$, $\varepsilon$ and $1$, and determinant $0$. So $\lambda_2+\lambda_3=2+2\varepsilon$ and $\lambda_2\lambda_3=1+\varepsilon+\varepsilon^2$, and by \eqref{eq:eigentime}
\[
\tau_0=\frac{\lambda_2+\lambda_3}{\lambda_2\lambda_3}=\frac{2(1+\varepsilon)}{1+\varepsilon+\varepsilon^2}<2.
\]

Now fix $t\in[0,2]$. Comparing the initial states $0$ and $2$ on the event $\{X_t=0\}$ gives
\begin{equation}\label{eq:d02}
\dbar_\varepsilon(t)\ge\Pp_0(X_t=0)-\Pp_2(X_t=0).
\end{equation}
The chain started at $0$ stays there for an exponential time with rate $\varepsilon$, so
\begin{equation}\label{eq:stay}
\Pp_0(X_t=0)\ge e^{-\varepsilon t}.
\end{equation}
For the chain started at $2$, let $Y_0=2,Y_1=1,Y_2,\dots$ be its jump chain and $S_0,S_1,\dots$ its holding times. Given the jump chain, the holding times are independent, and $S_m$ is exponential with rate $1$ if $Y_m=2$ and with rate $1+\varepsilon$ if $Y_m=1$. Let $N=\min\{m:Y_m=0\}$. Then $N\ge2$, $Y_{N-1}=1$ and
\[
T_0=S_0+\dots+S_{N-1}\ge S_0+S_{N-1}.
\]
Given the jump chain, $S_0$ and $S_{N-1}$ are independent and exponential with rates $1$ and $1+\varepsilon$; since these conditional laws do not depend on the jump chain, $S_0$ and $S_{N-1}$ are independent with these laws. Let $E$ and $E'$ be independent standard exponential variables and
\[
g(s)=\Pp(E+E'>s)=e^{-s}(1+s),\qquad s\ge0.
\]
Since $E+E'/(1+\varepsilon)\ge(E+E')/(1+\varepsilon)$,
\begin{equation}\label{eq:descend}
\Pp_2(X_t=0)\le\Pp_2(T_0\le t)\le\Pp\Bigl(E+\frac{E'}{1+\varepsilon}\le t\Bigr)\le\Pp\bigl(E+E'\le(1+\varepsilon)t\bigr)=1-g\bigl((1+\varepsilon)t\bigr).
\end{equation}
The function $g$ is nonincreasing and $(1+\varepsilon)t\le2+2\varepsilon$, so \eqref{eq:d02}, \eqref{eq:stay}, \eqref{eq:descend}, $e^{-\varepsilon t}\ge1-2\varepsilon$ and $e^{-2\varepsilon}\ge1-2\varepsilon$ give
\begin{equation}\label{eq:lower}
\dbar_\varepsilon(t)\ge e^{-\varepsilon t}-1+g\bigl((1+\varepsilon)t\bigr)\ge g(2+2\varepsilon)-2\varepsilon\ge e^{-2}(1-2\varepsilon)(3+2\varepsilon)-2\varepsilon.
\end{equation}
The right side of \eqref{eq:lower} exceeds $e^{-1}$ if and only if
\[
f(\varepsilon)=3-e-(4+2e^2)\varepsilon-4\varepsilon^2>0.
\]
The function $f$ is decreasing on $[0,\infty)$, and $e<2.7183$ and $e^2<7.3891$ give $f(1/100)>0.2817-0.1878-0.0004>0.09$. Hence, when $0<\varepsilon\le1/100$, $\dbar_\varepsilon(t)>e^{-1}$ for every $t\in[0,2]$, and $\tau_1>2$. This proves part~(i).

\subsection{Limits}
Let $k\ge1$ and $\varepsilon\to0$. The matrix $Q_\varepsilon$ converges to the matrix $Q_0$ with the same entries at $\varepsilon=0$, the generator of the pure death chain that moves from $i$ to $i-1$ at rate $1$ and is absorbed at $0$. The matrix $-Q_0$ is lower triangular with diagonal $(0,1,\dots,1)$, so its eigenvalues are $0$, which is simple, and $1$, with multiplicity $k$. The eigenvalues of a matrix depend continuously on its entries, and $0$ is an eigenvalue of every $-Q_\varepsilon$; so for small $\varepsilon$ the eigenvalue $0$ is the only eigenvalue of $-Q_\varepsilon$ near $0$, and the others, $\lambda_2,\dots,\lambda_{k+1}$, tend to $1$. By \eqref{eq:eigentime}, $\tau_0\to k$.

For $\tau_1$, let $P^\varepsilon_t=e^{tQ_\varepsilon}$ for $\varepsilon\ge0$. The identity $e^{tA}-e^{tB}=\int_0^te^{(t-s)A}(A-B)e^{sB}\,ds$ shows that $P^\varepsilon_t\to P^0_t$ uniformly for $t$ in compact sets, and hence $\dbar_\varepsilon\to\dbar_0$ uniformly on compact sets, where $\dbar_0(t)$ is defined from $P^0_t$ in the same way. Under $Q_0$ the chain started at $i$ passes through $i-1,\dots,1$ to $0$ with independent standard exponential holding times and then stays at $0$, so $P^0_t(0,\cdot)$ is the point mass at $0$ and $P^0_t(j,0)=\Pp(\Gamma_j\le t)$, where $\Gamma_0=0$. Let $i,j$ be states, $\mu=P^0_t(i,\cdot)$ and $\nu=P^0_t(j,\cdot)$, and suppose that $\mu(0)\ge\nu(0)$. Then
\[
\|\mu-\nu\|=\sum_x(\mu(x)-\nu(x))^+\le\mu(0)-\nu(0)+\sum_{x\ne0}\mu(x)=1-\nu(0)=\Pp(\Gamma_j>t)\le\Pp(\Gamma_k>t),
\]
with equality throughout when $i=0$ and $j=k$. Hence $\dbar_0(t)=\Pp(\Gamma_k>t)$. This function is continuous and strictly decreasing from $1$ to $0$. For $0<\delta<t_k$ it is therefore at least $\dbar_0(t_k-\delta)>e^{-1}$ on $[0,t_k-\delta]$, and $\dbar_0(t_k+\delta)<e^{-1}$. By uniform convergence on $[0,t_k+\delta]$, for all small $\varepsilon$ we have $\dbar_\varepsilon>e^{-1}$ on $[0,t_k-\delta]$ and $\dbar_\varepsilon(t_k+\delta)<e^{-1}$, so $t_k-\delta<\tau_1\le t_k+\delta$. This proves part~(ii).

\subsection{Discrete time}\label{sec:a-disc}
Let $0<c\le1/(1+\varepsilon)$. Every diagonal entry of $Q_\varepsilon$ is at least $-(1+\varepsilon)$, so $P=I+cQ_\varepsilon$ is a transition matrix, irreducible and reversible with respect to the same $\pi$. Its continuization has generator $P-I=cQ_\varepsilon$ and is $X^{(k,\varepsilon)}$ run at speed $c$, so its variation threshold, which is $\tau_1$ of $P$, is $\tau_1/c$, with $\tau_1$ that of $X^{(k,\varepsilon)}$. Continuization does not change mean hitting times \cite[Section~4.3.2]{AF}, so $\tau_0$ of $P$ is $\tau_0/c$. Hence $P$ has the same ratio $\tau_1/\tau_0$ as $X^{(k,\varepsilon)}$: for $k=2$ and $0<\varepsilon\le1/100$ it has $\tau_1>\tau_0$, and its ratio tends to $t_k/k$ as $\varepsilon\to0$.

For $k=2$ the inequality also holds, once $c$ is small, for the variation threshold of $P$ itself, $\tau_1^{\mathrm{disc}}=\min\{m\ge0:\dbar(m)\le e^{-1}\}$, where $\dbar(m)$ is computed from the matrix $P^m$ \cite[Section~4.3.1]{AF}. Let $0<\varepsilon\le1/100$ and $Q=Q_\varepsilon$, and write $\|A\|=\max_i\sum_j|A_{ij}|$. The matrices $I+cQ$ and $e^{cQ}$ are stochastic, so
\[
\|(I+cQ)^m-e^{mcQ}\|\le m\,\|I+cQ-e^{cQ}\|\le m\,\frac{(c\|Q\|)^2}{2}\,e^{c\|Q\|},
\]
which is at most $c\|Q\|^2e^{c\|Q\|}$ when $mc\le2$. The function $\dbar$ of a stochastic matrix changes by at most $\|A-B\|$ when the matrix $B$ is replaced by $A$, and by \eqref{eq:lower} and the bound $f(\varepsilon)\ge f(1/100)>0.09$, $\dbar_\varepsilon(t)-e^{-1}\ge e^{-2}f(\varepsilon)>0.01$ for $t\in[0,2]$. Hence, when $c\|Q\|^2e^{c\|Q\|}\le0.01$, the matrix $P^m$ has $\dbar>e^{-1}$ for every $m\le2/c$, and $\tau_1^{\mathrm{disc}}>2/c>\tau_0/c$.

\section{Proof of Theorem~\ref{thm:b}}\label{sec:b}

\begin{proof}[Proof of Theorem~\ref{thm:b}]
Replacing $u$ by $-u$ we may assume $u(i)\ge0$. Since $\pi Q=0$, we have $-\lambda_2\pi(u)=\pi(Qu)=0$, so $\pi(u)=0$. Let $U$ be a stopping time for a filtration $(\mathcal F_t)$ with respect to which $X$ is Markov, with $\Pp_i(X_U\in\cdot)=\pi$, and assume that $\E_iU<\infty$. For every function $h$ on $I$, the process $h(X_t)-h(X_0)-\int_0^t(Qh)(X_s)\,ds$ is an $(\mathcal F_t)$-martingale, so optional stopping at the bounded time $U\wedge t$, with $h=u$, gives
\[
\E_iu(X_{U\wedge t})-u(i)=-\lambda_2\,\E_i\int_0^{U\wedge t}u(X_s)\,ds.
\]
Let $t\to\infty$. Since $U<\infty$ almost surely and $|u|\le1$, the left side tends to $\E_iu(X_U)-u(i)=\pi(u)-u(i)=-u(i)$ by bounded convergence, and the right side tends to $-\lambda_2\E_i\int_0^Uu(X_s)\,ds$ by dominated convergence, since the integral is at most $U$ in absolute value. Therefore
\[
u(i)=\lambda_2\,\E_i\int_0^Uu(X_s)\,ds\le\lambda_2\,\E_iU,
\]
that is, $\E_iU\ge u(i)/\lambda_2=\tau_2|u(i)|$. Taking $i$ with $|u(i)|=1$ gives $\inf_U\E_iU\ge\tau_2$, and so $\tau_1^{(2)}\ge\tau_2$.
\end{proof}

\subsection*{Discrete time}
Let $P$ be an irreducible reversible transition matrix, $\lambda_2<1$ its second largest eigenvalue, and $u$ a real function with $Pu=\lambda_2u$, $\max_k|u(k)|=1$ and $u(i)\ge0$. Then $\pi(u)=\pi(Pu)=\lambda_2\pi(u)$, so $\pi(u)=0$, and $u(X_n)-u(X_0)+(1-\lambda_2)\sum_{m<n}u(X_m)$ is a martingale. The argument above, with sums in place of integrals, gives
\[
u(i)=(1-\lambda_2)\,\E_i\sum_{m<U}u(X_m)\le(1-\lambda_2)\,\E_iU,
\]
so $\E_iU\ge\tau_2|u(i)|$ with $\tau_2=1/(1-\lambda_2)$, and $\tau_2\le\tau_1^{(2)}$.

\subsection*{Remarks}
1. For the chain on $\{1,2\}$ that jumps from $1$ to $2$ at rate $a$ and from $2$ to $1$ at rate $b$, $-Q$ has eigenvalues $0$ and $a+b$, so $\tau_2=1/(a+b)$, and $\pi=(b,a)/(a+b)$. Started at $1$, the stopping time that equals $0$ with probability $b/(a+b)$ and $T_2$ otherwise, the choice being independent of the chain, has $X_U$ distributed as $\pi$ and mean $\frac a{a+b}\cdot\frac1a=\frac1{a+b}$, and similarly from $2$. With Theorem~\ref{thm:b} this gives $\tau_1^{(2)}=1/(a+b)=\tau_2$. For this chain also $\dbar(t)=e^{-(a+b)t}$ and $\tau_1=\tau_0=1/(a+b)$.

2. Started at $i$, run the chain until it hits a random state with law $\pi$ chosen independently of the chain. This stopping time has $X_U$ distributed as $\pi$ and mean $\sum_j\pi_j\E_iT_j=\tau_0$, so $\tau_1^{(2)}\le\tau_0$ \cite[Section~4.3.2]{AF}. With Theorem~\ref{thm:b}, $\tau_2\le\tau_1^{(2)}\le\tau_0$, a refinement of the inequality $\tau_2\le\tau_0$ of \cite[Lemma~4.24]{AF}.

3. Let $Z_{ij}=\int_0^\infty(\Pp_i(X_t=j)-\pi_j)\,dt$. By \cite[(4.47)]{AF}, $\tau_1^{(2)}=\max_{i,j}(-Z_{ij}/\pi_j)$, so Theorem~\ref{thm:b} says that $\max_{i,j}(-Z_{ij}/\pi_j)\ge1/\lambda_2$. This can also be seen directly. Take $u$ as in Theorem~\ref{thm:b}, replaced by $-u$ if necessary so that $u(i)=1$ for some state $i$. Since $\pi(u)=0$, $\sum_jZ_{ij}u(j)=\int_0^\infty e^{-\lambda_2t}u(i)\,dt=\tau_2$. Let $m=\max_j(-Z_{ij}/\pi_j)$, so that $Z_{ij}+m\pi_j\ge0$ for every $j$. Since $\sum_jZ_{ij}=0$ and $\pi(u)=0$,
\[
\tau_2=\sum_j(Z_{ij}+m\pi_j)u(j)\le\sum_j(Z_{ij}+m\pi_j)=m.
\]

4. The proof of Theorem~\ref{thm:b} uses reversibility only to know that $\lambda_2$ is real with a real eigenfunction. For any finite irreducible chain in continuous time, any eigenvalue $\lambda\ne0$ of $-Q$ with eigenvector $u\in\mathbb C^I$, and any probability measure $\mu$ on $I$, the same argument applied to the real and imaginary parts of $u$ shows that every randomized stopping time $U$ with $\Pp_i(X_U\in\cdot)=\mu$ satisfies $|\lambda|\max_k|u(k)|\,\E_iU\ge|u(i)-\mu(u)|$. For $\mu=\pi$ this gives $\tau_1^{(2)}\ge1/|\lambda|$.

5. In discrete time Theorem~\ref{thm:b} does not extend to the absolute relaxation time $1/(1-\max(\lambda_2,|\lambda_n|))$, where $\lambda_n$ is the smallest eigenvalue of $P$. For
\[
P=\begin{pmatrix}\delta&1-\delta\\1-\delta&\delta\end{pmatrix},\qquad0<\delta<\tfrac12,
\]
$\lambda_2=\lambda_n=2\delta-1$ and $\tau_2=1/(2(1-\delta))$. Stopping at time $0$ with probability $1/2$ and otherwise at the first visit to the other state gives a stopping time with $X_U$ distributed as $\pi$ and mean $1/(2(1-\delta))$, so $\tau_1^{(2)}=1/(2(1-\delta))$ by Theorem~\ref{thm:b}, while the absolute relaxation time is $1/(2\delta)$.

\section{Computations}\label{sec:checks}

The repository listed under Data availability contains a Python script, \texttt{checks.py}, with its output, and a Lean file, \texttt{Margin.lean}. For $X^{(2,\varepsilon)}$ with $\varepsilon=1/100$ and $\varepsilon=1/1000$, the script computes $\tau_0$ from the closed form, from the eigentime identity and from the mean hitting times; computes $\dbar_\varepsilon(t)$ from the matrix exponential and from the spectral decomposition, in $40$-digit arithmetic; checks $\dbar_\varepsilon(t)>e^{-1}$, \eqref{eq:d02}, \eqref{eq:stay}, the first inequality of \eqref{eq:lower} and the bound $e^{-\varepsilon t}-1+g((1+\varepsilon)t)>e^{-1}$ on a grid of step $0.001$ on $[0,2]$, and \eqref{eq:descend} at the multiples of $0.05$ in $[0,2]$; and finds $\tau_1$ by bisection. At $\varepsilon=1/100$, $\tau_0=1.99980\ldots$, $\tau_1=2.13306\ldots$ and the smallest value of $\dbar_\varepsilon$ on $[0,2]$ is $0.40240\ldots$; at $\varepsilon=1/1000$, $\tau_1/\tau_0=1.07244\ldots$. At $\varepsilon=1/20$, where $f(\varepsilon)<0$ and the proof does not apply, the ratio $\tau_1/\tau_0$ is still $1.0421\ldots$. For the chains on $\{0,\dots,k\}$ the script computes $t_k$ for $k\le60$ by bisection, with a sign check on an interval of width $2\cdot10^{-40}$ for $k\le8$; checks $\dbar_0(t)=\Pp(\Gamma_k>t)$, attained at the pair $(0,k)$, for $k=2,3,4,6$; and computes $\tau_1/\tau_0$ at $\varepsilon=10^{-3},10^{-5},10^{-7}$ for $k=2,3,4,6$, where it approaches $t_k/k$. For Theorem~\ref{thm:b} it draws $3000$ random reversible chains with $2$ to $8$ states in continuous time and $3000$ in discrete time. For each it computes $\tau_1^{(2)}$ from \cite[(4.47)]{AF} and checks the inequality of Theorem~\ref{thm:b} at every state. For $600$ of the chains in each sample it also solves, at every state $i$, the linear program that minimizes the total mass of a measure $x\ge0$ on $I$ with $xG=\pi-\delta_i$, where $G=Q$ or $G=P-I$ (the expected occupation measure of a stopping time $U$ with $\Pp_i(X_U\in\cdot)=\pi$ satisfies these constraints and has total mass $\E_iU$), checks that its value agrees with \cite[(4.47)]{AF}, and checks the mean hitting times against the fundamental matrix. The smallest value of $\tau_1^{(2)}/\tau_2$ in each sample is $1.000000$ to six decimals. The script also checks Remark~1 and the example of Remark~5; in the discrete-time sample, the absolute relaxation time exceeds $\tau_1^{(2)}$ for $1584$ of the $3000$ chains. The Lean file, checked with Lean~4 \cite{MU21} (version 4.33.1) and Mathlib \cite{Mathlib20}, proves that $g$ is nonincreasing, that $f(\varepsilon)>0$ for $0<\varepsilon\le1/100$, and that $e^{-\varepsilon t}-1+g((1+\varepsilon)t)>e^{-1}$ for $0<\varepsilon\le1/100$ and $0\le t\le2$, which is the real-variable part of the proof of Theorem~\ref{thm:a}(i).

These computations test the identities and inequalities used in the proofs; they are not part of them.

\section*{Acknowledgements and funding}
The research presented in this paper was supported by the European Research
Council (ERC) under the European Union's Horizon Europe research and
innovation programme (grant agreement No.~101041711), by the Simons
Foundation as part of the Collaboration on the Mathematical and Scientific
Foundations of Deep Learning, by Heights Labs, by Convex Nexus Capital, by
the Israel Science Foundation (grant number 2258/19), and by the Israel
Science Foundation (ISF Grant 4101/25). Computational resources were
provided in part by Google Cloud through its education credits program
(credit 530550103).

\section*{Declaration of competing interest}
The author holds equity in Heights Labs and in Convex Nexus Capital and
has a management role at Convex Nexus Capital. These entities are included
in the funding declaration above.

\section*{Declaration of generative AI and AI-assisted technologies}
The paper comes out of Math Brain, an ongoing project on the use of AI systems for research mathematics. Open Problem~4.48 of \cite{AF} was among the open problems in probability on which the project's agents worked in September 2026. On 6 September 2026 an agent running in OpenAI's Codex wrote the optional stopping argument that proves Theorem~\ref{thm:b}. On 14 September 2026 Claude Fable~5.1 (Anthropic) found the chains of Theorem~\ref{thm:a}, sketched the proof of its part~(ii), and wrote out the proof of Theorem~\ref{thm:b}. On 13 October 2026 Claude Opus~5.5 (Anthropic), working in Claude Code, proved part~(i) of Theorem~\ref{thm:a} for $0<\varepsilon\le1/100$, completed the proof of part~(ii), searched the literature, and wrote the remarks, the computer checks, the Lean file and the manuscript. The author directed the work and is responsible for the content of the paper and for decisions concerning its dissemination.

\section*{Data availability}
No data were used. The code and the Lean file are available at\\ \url{https://github.com/yspennstate/variation-threshold-stationary-times}.

\begin{thebibliography}{9}

\bibitem{AF}
D.~Aldous and J.~A.~Fill, \emph{Reversible Markov Chains and Random Walks on Graphs}, unfinished monograph, 2002, recompiled version 2014, \url{https://www.stat.berkeley.edu/~aldous/RWG/book.pdf}.

\bibitem{MU21}
L.~de~Moura and S.~Ullrich, The Lean~4 theorem prover and programming language, in \emph{Automated Deduction -- CADE~28}, Lecture Notes in Comput. Sci. \textbf{12699}, Springer, 2021, 625--635.

\bibitem{Mathlib20}
The mathlib Community, The Lean mathematical library, in \emph{Proceedings of the 9th ACM SIGPLAN International Conference on Certified Programs and Proofs (CPP 2020)}, ACM, 2020, 367--381.

\end{thebibliography}
\end{document}
